% Generated draft: requires mathematical and source review.
\documentclass[11pt]{article}
\usepackage[margin=25mm]{geometry}
\usepackage{fontspec}
\setmainfont{Latin Modern Roman}
\usepackage{amsmath,amssymb,mathtools,unicode-math}
\setmathfont{Latin Modern Math}
\usepackage{enumitem,needspace,xurl,hyperref}
\setlength{\parindent}{0pt}
\setlength{\parskip}{5pt}
\setlength{\emergencystretch}{3em}
\newcommand{\N}{\mathbb{N}}
\newcommand{\Z}{\mathbb{Z}}
\newcommand{\Q}{\mathbb{Q}}
\newcommand{\R}{\mathbb{R}}
\newcommand{\C}{\mathbb{C}}
\newcommand{\F}{\mathbb{F}}
\newcommand{\A}{\mathbb{A}}
\newcommand{\PP}{\mathbb{P}}
\newcommand{\E}{\mathbb{E}}
\newcommand{\eps}{\varepsilon}
\newcommand{\dd}{\,\mathrm{d}}
\begin{document}
\section*{UnsolvedMath ID: 20001410}
\textbf{Source problem code:} catalogue entry
\subsection{Definitions and mathematical statement}
Problem 25 asks for hitting-time and extreme-value laws in billiards when a shrinking phase-space target is a strip induced by a configuration-space set rather than a phase-space ball. The source does not identify the billiard class, target geometry, collision versus flow observation, invariant measure or singularity assumptions. The appended ``Open Problems and Questions: Thursday'' is the next heading, not part of the target definition.

The aid treats planar collision strips. For a bounded planar billiard configuration domain \(Q\), let \(L\) be its total boundary length, \(s\) boundary arclength, and \(\phi\in(-\pi/2,\pi/2)\) the outgoing angle from the inward normal. With unit-speed collision dynamics \(T\), normalize invariant collision measure by
\[
\dd\mu(s,\phi)=\frac{\cos\phi}{2L}\,\dd s\,\dd\phi.
\]
A full-angle strip \(A_I=\{(s,\phi):s\in I\}\) over a boundary arc has exact mass \(\mu(A_I)=|I|/L\). A two-sided arclength strip \(|s-s_0|<\varepsilon\), away from boundary-coordinate endpoints, therefore has mass \(2\varepsilon/L\).

For a compact strictly convex target \(K\) lying inside the free configuration domain, let \(B_K\) be the collision states whose next free-flight segment enters \(K\). In the flux setting described by the aid, transporting Liouville flux to the target boundary gives
\[
\mu(B_K)=\frac{\operatorname{Per}(K)}L.
\]
Convexity ensures that a segment does not make multiple separate entries into \(K\); the collision/flux correspondence must be well defined up to null grazing or singular states. For homothetic targets \(K_\varepsilon=q_0+\varepsilon K_0\) contained in the free domain, \(\mu(B_{K_\varepsilon})=\varepsilon\operatorname{Per}(K_0)/L\). These are collision-event masses, not the configuration-area mass or a physical-time entry rate.

Conditionally, if a target family \(A_\varepsilon\) of mass \(a_\varepsilon\sim C\varepsilon\), \(C>0\), has an extremal-index-one Poisson visit law on the collision scale, then for \(\tau_{A_\varepsilon}=\inf\{j\ge1:T^jx\in A_\varepsilon\}\),
\[
\mu\{a_\varepsilon\tau_{A_\varepsilon}>t\}\longrightarrow e^{-t}\quad(t\ge0).
\]
For an observable whose exceedance set at level \(u\) is \(A_{e^{-u}}\), the associated maximum \(M_n=\max_{0\le j<n}\psi\circ T^j\) has the Gumbel normalization
\[
u_n(y)=\log(Cn)+y,\qquad\mu\{M_n\le u_n(y)\}\longrightarrow\exp(-e^{-y}),
\]
under the same Poisson hypothesis. The scale used in \(e^{-u}\) fixes the arclength reference units. These limit laws are conditional conclusions, not consequences of the mass calculation alone.

The aid's fixed-lag clustering diagnostic is restricted to finitely many uniformly regular branches. For fixed \(k\), restrict the returning boundary coordinate \(s_k(s,\phi)\) to the central fiber \(s=s_0\). If its returns to \(s_0\) are isolated of finite contact order \(m\), with a uniform local bound \(|s_k(s_0,\phi)-s_0|\ge c|\phi-\phi_0|^m\), bounded derivatives in the transverse coordinate, and the required finite-branch control, then
\[
\mu(A_\varepsilon\cap T^{-k}A_\varepsilon)=O(\varepsilon^{1+1/m})=o(\mu(A_\varepsilon)).
\]
For several contact orders, the estimate is taken branchwise. In contrast, an interval of persistent fiber overlap, with regularity and positive-density control, has intersection of order \(\varepsilon\) and can cause clustering. Isolated finite-order contact and persistent overlap are therefore distinct mechanisms. This fixed-lag diagnostic does not control arbitrarily many branches, grazing singularities, growing lags or all the mixing/short-return requirements for a Poisson theorem; persistent overlap alone is not a complete formula for the extremal index.

The aid reports that the boundary-strip version is substantially answered by Bunimovich and Su's 2024 Poisson theorem for billiard-boundary holes. Their precise billiard hypotheses and theorem citation are not supplied. The registry labels the entry partially solved; attempt 1 is a partial result, problem partially solved at the run, internal rigor proved and novelty confidence low. The machine-generated aid is not peer reviewed and requires independent verification before citation.
\subsection{Short English statement}
The question concerns strip targets determined by position rather than phase-space balls. The aid gives exact collision-measure normalizations and conditional exponential/Gumbel laws, then distinguishes negligible isolated finite-order returns from persistent fiber overlap that can create clustering. These calculations do not replace the regularity and mixing hypotheses of a billiard Poisson theorem.
\subsection{Sources}
\begin{itemize}
\item Local UnsolvedMath registry, permanent ID \(20001410\), source code \texttt{catalogue entry}; \texttt{aim-dynamical-systems-notes.json}, \texttt{notes[67]}.
\item AIM, \emph{Stochastic methods for non-equilibrium dynamical systems}, Problem 25, with the ensuing Thursday heading: \url{https://aimath.org/pastworkshops/noneqdynsysproblems.pdf}.
\item Bunimovich and Su (2024), Poisson theorem for billiard-boundary holes, as reported in the research summary; no article title, publication, theorem number or exact model hypotheses supplied. The original literature field is empty.
\item Supplied research attempt 1 and registry update August 14, 2026; no independent source visit or verification of the literature status is claimed.
\end{itemize}
\end{document}
